\SourceLeaf{63}
\begin{center}\emph{Methodus radices primitivas assignandi.}\end{center}

\Article{73}

Methodi radices primitivas inveniendi maximam partem tentando innituntur. Si quis ea quae art. 55 docuimus cum iis quae infra de solutione congruentiae $x^n\equiv1$ trademus confert, omnia fere, quae per methodos directas effici possunt, habebit. Ill. Euler confitetur, \emph{Opusc. Analyt. T.} I. \emph{p.} 152, maxime difficile videri, hos numeros assignare, eorumque indolem ad profundissima numerorum mysteria esse referendam. At tentando satis expedite sequenti modo determinari possunt. Exercitatus operationis prolixitati per multifaria artificia particularia succurrere sciet: haec vero per usum multo citius quam per praecepta ediscuntur.

$1^0$. Assumatur ad libitum numerus ad $p$ (ita semper modulum designamus) primus, $a$, (plerumque ad calculi brevitatem conducit, si quam minimum accipimus, \emph{ex. gr.} numerum $2$) determineturque eius periodus (art. 46), \emph{i. e.} residua minima ipsius potestatum, donec ad potestatem $a^t$ perveniatur, cuius residuum minimum sit $1$\HistoricalNote{*)}{Quisquis sponte perspiciet, non opus esse has potestates ipsas novisse, quum cuiusvis residuum minimum facile ex residuo minimo potestatis praecedentis obtineri possit.}. Iam si fuerit $t=p-1$, $a$ est radix primitiva.

$2^0$. Si vero $t<p-1$, accipiatur alius numerus $b$ in periodo ipsius $a$ non contentus, investigeturque simili modo huius periodus. Designato exponente ad quem $b$ pertinet per $u$, facile perspicitur $u$ neque ipsi $t$ aequalem neque ipsius partem aliquotam esse posse, in utroque enim casu fieret $b^t\equiv1$, quod esse nequit, quum periodus ipsius $a$ omnes numeros amplectatur, quorum potestas exponentis $t$ unitati congrua (art. 53). Quodsi $u$ fuerit $=p-1$, erit $b$ radix primitiva; si vero $u$ non quidem $=p-1$, sed tamen multiplum ipsius $t$, id lucrati sumus, ut numerus constet ad exponentem maiorem pertinens, adeoque scopo nostro, qui est invenire numerum ad exponentem \emph{maximum} pertinentem, propiores iam simus. Si vero $u$ neque $=p-1$, neque ipsius $t$ multiplum, tamen numerum invenire possumus ad exponentem ipsis $t,u$ maiorem pertinentem, nempe ad exponentem minimo dividuo communi numerorum $t,u$ aequalem. Sit hic $=y$, resolvaturque $y$ ita in duos factores inter se primos, $m,n$, ut alter ipsum $t$, alter ipsum $u$ metiatur\HistoricalNote{\ensuremath{\dagger})}{Quomodo hoc fieri possit ex art. 18 haud difficulter derivatur. Resolvatur $y$ in factores tales, qui sint aut numeri primi diversi aut numerorum primorum diversorum potestates. Horum quisque alterutrum numerorum $t,u$ metietur (sive etiam utrumque). Adscribantur singuli aut numero $t$ aut numero $u$, prout illum aut hunc metiuntur: quando aliquis utrumque metitur, arbitrarium est, cui adscribatur: productum ex iis qui ipsi $t$ adscripti sunt, sit $=m$, productum e reliquis $=n$, facileque perspicietur $m$ ipsum $t$, $n$ ipsum $u$ metiri, atque esse $mn=y$.}. Tum fiat potestas $\frac{t}{m}^{\mathrm{ta}}$ ipsius $a$, $\equiv A$, pote\SourceLeaf{64}stas $\frac{u}{n}^{\mathrm{ta}}$ ipsius $b$, $\equiv B\,(\operatorname{mod.}p)$, eritque productum $AB$ numerus ad exponentem $n$ pertinens; facile enim intelligitur, $A$ ad exponentem $m$, $B$ ad exponentem $n$ pertinere; adeoque productum $AB$ ad $mn$ pertinebit, quia $m,n$ inter se sunt primi, id quod prorsus eodem modo uti in art. 55, II processimus probari poterit.

$3^0$. Iam si $y=p-1$, $AB$ erit radix primitiva; sin minus, simili modo ut antea alius numerus adhibendus erit, in periodo ipsius $AB$ non occurrens; eritque hic aut radix primitiva, aut pertinebit ad exponentem ipso $y$ maiorem. aut certe ipsius auxilio (uti ante) numerus ad exponentem ipso $y$ maiorem pertinens inveniri poterit. Quum igitur numeri qui per repetitionem huius operationis prodeunt, ad exponentes continuo crescentes pertineant, manifestum est tandem numerum inventum iri, qui ad exponentem \emph{maximum} pertineat, \emph{i. e.} radicem primam, \emph{q. e. f.}

\Article{74}

Per exemplum praecepta haec clariora fient. Sit $p=73$, pro quo radix primitiva quaeratur. Tentemus primo numerum $2$, cuius periodus prodit haec:

\[\setlength{\arraycolsep}{2.5pt}\begin{array}{@{}rrrrrrrrrrl@{}}
1.&2.&4.&8.&16.&32.&64.&55.&37.&1&\text{etc.}\\
0.&1.&2.&3.&4.&5.&6.&7.&8.&9&\text{etc.}
\end{array}\]

Quum igitur iam potestas exponentis $9$ unitati congrua fiat, $2$ non est radix primitiva. Tentetur alius numerus in periodo ipsius $2$ non occurrens \emph{ex. gr.} $3$, cuius periodus est haec:

\[\setlength{\arraycolsep}{2.5pt}\begin{array}{@{}rrrrrrrrrrrrrl@{}}
1.&3.&9.&27.&8.&24.&72.&70.&64.&46.&65.&49.&1&\text{etc.}\\
0.&1.&2.&3.&4.&5.&6.&7.&8.&9.&10.&11.&12&\text{etc.}
\end{array}\]

Quare neque $3$ est radix primitiva. Exponentium autem ad quos $2,3$ pertinent, (\emph{i. e.} numerorum $9,12$) dividuus communis minimus est $36$, qui in factores $9$ et $4$ ad praecepta art. praec. resolvitur. Evehendus itaque $2$ ad potestatem exponentis $\frac{9}{9}$, \emph{i. e.} numerus $2$ ipse retinendus; $3$ autem ad potestatem exponentis $3$: productum ex his est $54$, quod itaque ad exponentem $36$ pertinebit. Si denique ipsius $54$ periodus computatur numerusque in hac non contentus \emph{ex. gr.} $5$ denuo tentatur, hunc esse radicem primitivam, reperietur.

\SourceLeaf{65}
\begin{center}\emph{Theoremata varia de periodis et radicibus primitivis.}\end{center}

\Article{75}

Antequam hoc argumentum deseramus, propositiones quasdam trademus, quae ob simplicitatem suam attentione haud indignae videntur.

\emph{Productum ex omnibus terminis periodi numeri cuiusvis est $\equiv1$, quando ipsorum multitudo, sive exponens ad quem numerus pertinet, est impar, et $\equiv-1$, quando ille exponens est par.}

\emph{Ex.} Pro modulo $13$ periodus numeri $5$ constat ex his terminis $1,5,12,8$ quorum productum $480\equiv-1\,(\operatorname{mod.}13)$.

Secundum eundem modulum periodus numeri $3$ constat e terminis $1,3,9$ quorum productum $27\equiv1\,(\operatorname{mod.}13)$.

\emph{Demonstr.} Sit exponens, ad quem numerus pertinet, $t$, atque index numeri, $\frac{p-1}{t}$, id quod si basis rite determinatur semper fieri potest (art. 71). Tum index producti ex omnibus periodi terminis erit

\[\equiv(1+2+3+\text{etc.}+t-1)\frac{p-1}{t}=\frac{(t-1)(p-1)}{2}\]

\emph{i. e.} $\equiv0\,(\operatorname{mod.}p-1)$, quando $t$ impar, et $\equiv\frac{p-1}{2}$, quando $t$ par; hinc in priori casu productum illud $\equiv1\,(\operatorname{mod.}p)$; in posteriori vero $\equiv-1\,(\operatorname{mod.}p)$, (art. 62). \emph{Q. E. D.}

\Article{76}

Si numerus iste in theor. praecedente est radix primitiva, eius periodus omnes numeros $1,2,3\mathpunct{.\,.\,.\,.\,.\,.}\,p-1$ comprehendet, quorum productum itaque semper $\equiv-1$ (namque $p-1$ semper par, unico casu $p=2$ excepto in quo $-1$ et $+1$ aequivalent). Theorema hoc elegans quod ita enunciari solet: \emph{productum ex omnibus numeris numero primo dato minoribus, unitate auctum per hunc primum est divisibile}, primum a cel. Waring est prolatum armigeroque Wilson adscriptum, \emph{Meditt. algebr. Ed.} 3. \emph{p.} 380. Sed neuter demonstrare potuit, et cel. Waring fatetur demonstrationem eo difficiliorem videri, quod nulla \emph{notatio} fingi possit, quae numerum primum exprimat. --- At nostro quidem iudicio huiusmodi veritates ex notionibus potius quam ex notationibus hauriri debebant. Postea ill. La Grange demonstrationem dedit, \emph{Nouv. M\'em. de l'Ac. de Berlin,} 1771. Innititur ea consi\SourceLeaf{66}derationi coefficientium ex evolutione producti

\[x+1.\,x+2.\,x+3\mathpunct{.\,.\,.\,.}\,x+p-1\]

\noindent oriundorum. Scilicet posito hoc producto

\[\equiv x^{p-1}+Ax^{p-2}+Bx^{p-3}+\text{etc.}+Mx+N\]

\noindent coefficientes $A,B$ etc. $M$ per $p$ erunt divisibiles, $N$ vero erit $=1.2.3\mathpunct{.\,.\,.\,.}\,p-1$. Iam pro $x=1$, productum per $p$ divisibile; tunc autem erit $\equiv1+N\,(\operatorname{mod.}p)$; quare necessario $1+N$ per $p$ dividi poterit.

Denique ill. Euler in \emph{Opusc. analyt. T.} I. \emph{p.} 329 demonstrationem dedit, cum ea quam nos hic exposuimus conspirantem. Quodsi tales viri theorema hoc meditationibus suis non indignum censuerunt, non improbatum iri speramus, si aliam adhuc demonstrationem apponimus.

\Article{77}

Quando secundum modulum $p$, productum duorum numerorum $a,b$ unitati est congruum, numeros $a,b$ cum ill. Euler \emph{socios} vocemus. Tum secundum sect. praec. quivis numerus positivus ipso $p$ minor socium habebit positivum ipso $p$ minorem et quidem unicum. Facile autem probari potest ex numeris $1,2,3\mathpunct{.\,.\,.\,.}\,p-1$; $1$ et $p-1$ esse unicos qui sibi ipsis sint socii: numeri enim sibi ipsis socii, radices erunt congruentiae $xx\equiv1$; quae quoniam est secundi gradus plures quam duas radices, \emph{i. e.} alias quam $1$ et $p-1$ habere nequit. Abiectis itaque his numerorum reliquorum $2.3\mathpunct{.\,.\,.\,.}\,p-2$ bini semper erunt associati; quare productum ex ipsis erit $\equiv1$ adeoque productum ex omnibus $1,2,3\mathpunct{.\,.\,.}\,p-1$, $\equiv p-1$ sive $\equiv-1$. \emph{Q. E. D.}

\emph{Ex. gr.} pro $p=13$ numeri $2,3,4\mathpunct{.\,.\,.\,.}\,11$ ita associantur: $2$ cum $7$; $3$ cum $9$; $4$ cum $10$; $5$ cum $8$; $6$ cum $11$; scilicet $2.7\equiv1$; $3.9\equiv1$ etc. Hinc $2.3.4\mathpunct{.\,.\,.\,.}\,11\equiv1$; adeoque $1.2.3\mathpunct{.\,.\,.\,.}\,12\equiv-1$.

\Article{78}

Potest autem theorema Wilsonianum generalius sic proponi. \emph{Productum ex omnibus numeris, numero quocunque dato $A$ minoribus simulque ad ipsum primis, congruum est secundum $A$, unitati vel negative vel positive sumtae.} Negative sumenda est unitas, quando $A$ est formae $p^m$. aut huiusce $2p^m$, designante $p$ nume%
